\documentclass[withoutpreface,bwprint]{cumcmthesis} %去掉封面与编号页，电子版提交的时候使用。


\usepackage[framemethod=TikZ]{mdframed}
\usepackage{url}   % 网页链接
\usepackage{subcaption} % 子标题
\title{基于Arima模型对北京市房山区气温的时间序列预测}
\tihao{}
\baominghao{}
\schoolname{北京理工大学}
\membera{黄宇辰}
\memberb{彭子瀚}
\memberc{ }
\supervisor{ }%辅导老师
\yearinput{2025}
\monthinput{9}
\dayinput{1}

\begin{document}

 \maketitle
 \begin{abstract}


\keywords{\TeX{}\quad  图片\quad   表格\quad  公式}
\end{abstract}


\section{引言}
本文是北京理工大学2025年9月8日完成的数学建模课程大作业, 题目为"基于Arima模型对北京市房山区气温的时间序列预测". 本文旨在通过对北京市房山区气温数据的分析与建模, 利用Arima模型对未来气温进行预测, 以期为相关领域提供参考依据. 文章结构如下: 第二部分介绍问题背景与现有研究; 第三部分阐述模型假设与符号说明; 第四部分详细描述模型的建立与求解过程; 第五部分进行模型的检验与结果分析,并与实际数据进行对比,讨论产生差异的可能原因; 最后总结全文,并提出未来研究的方向.
\section{问题背景分析}
北京市位于我国华北平原北端，地势整体呈现出西北高、东南低的特征，总面积16410.54 km2，山区面积占62\%，平原区面积占38\%，平均海拔43.5 m，森林覆盖率44.4\%。北京市气候为典型的温带季风性气候，春季3～5月、夏季6～8月、秋季9～11月、冬季12～2月四季分明，冬半年受冷涡系统影响多偏北风，夏半年受东亚季风影响多偏南风。近十多年来，北京市平均气温12℃，月均气温变化范围为-4℃（1月）～26℃（7月）；北京市年累计降水量550 mm左右。数十年来，北京市经济持续快速发展，城市建成区面积不断扩大。(\cite{cheng2025})



\subsection{现有研究的明显短板}


\section{模型假设与符号说明}




\section{模型的建立与求解}

\section{模型的检验}


%参考文献
\begin{thebibliography}{9}%宽度9
   \bibitem{cheng2025beijing}
程兵芬, 胡晓征, 张瑞, 等.
\newblock 北京地区降水量时空变异特征分析
\newblock \emph{华北科技学院学报} (或实际期刊名), 2025.
\newblock DOI: 10.19956/j.cnki.ncist.2025.04.014.
\end{thebibliography}

\newpage
%附录
\begin{appendices}

\section{模板所用的宏包}
\begin{table}[htbp]
    \centering
    \caption{宏包罗列}
    \begin{tabular}{ccccc}
        \toprule
        \multicolumn{5}{c}{模板中已经加载的宏包} \\
        \midrule
        amsbsy & amsfonts & {amsgen} & {amsmath} & {amsopn} \\
        amssymb & amstext & {appendix} & {array} & {atbegshi} \\
        atveryend & auxhook & {bigdelim} & {bigintcalc} & {bigstrut} \\
        bitset & bm    & {booktabs} & {calc} & {caption} \\
        caption3 & CJKfntef & {cprotect} & {ctex} & {ctexhook} \\
        ctexpatch & enumitem & {etexcmds} & {etoolbox} & {everysel} \\
        expl3 & fix-cm & {fontenc} & {fontspec} & {fontspec-xetex} \\
        geometry & gettitlestring & {graphics} & {graphicx} & {hobsub} \\
        hobsub-generic & hobsub-hyperref & {hopatch} & {hxetex} & {hycolor} \\
        hyperref & ifluatex & {ifpdf} & {ifthen} & {ifvtex} \\
        ifxetex & indentfirst & {infwarerr} & {intcalc} & {keyval} \\
        kvdefinekeys & kvoptions & {kvsetkeys} & {l3keys2e} & {letltxmacro} \\
        listings & longtable & {lstmisc} & {ltcaption} & {ltxcmds} \\
        multirow & nameref & {pdfescape} & {pdftexcmds} & {refcount} \\
        rerunfilecheck & stringenc & {suffix} & {titletoc} & {tocloft} \\
        trig  & ulem  & {uniquecounter} & {url} & {xcolor} \\
        xcolor-patch & xeCJK & {xeCJKfntef} & {xeCJK-listings} & {xparse} \\
        xtemplate & zhnumber &       &       &  \\
        \bottomrule
    \end{tabular}%
    \label{tab:addlabel}%
\end{table}%

以上宏包都已经加载过了，不要重复加载它们。

\section{排队算法--matlab 源程序}

\begin{lstlisting}[language=matlab]
kk=2;[mdd,ndd]=size(dd);
while ~isempty(V)
[tmpd,j]=min(W(i,V));tmpj=V(j);
for k=2:ndd
[tmp1,jj]=min(dd(1,k)+W(dd(2,k),V));
tmp2=V(jj);tt(k-1,:)=[tmp1,tmp2,jj];
end
tmp=[tmpd,tmpj,j;tt];[tmp3,tmp4]=min(tmp(:,1));
if tmp3==tmpd, ss(1:2,kk)=[i;tmp(tmp4,2)];
else,tmp5=find(ss(:,tmp4)~=0);tmp6=length(tmp5);
if dd(2,tmp4)==ss(tmp6,tmp4)
ss(1:tmp6+1,kk)=[ss(tmp5,tmp4);tmp(tmp4,2)];
else, ss(1:3,kk)=[i;dd(2,tmp4);tmp(tmp4,2)];
end;end
dd=[dd,[tmp3;tmp(tmp4,2)]];V(tmp(tmp4,3))=[];
[mdd,ndd]=size(dd);kk=kk+1;
end; S=ss; D=dd(1,:);
 \end{lstlisting}

 \section{规划解决程序--lingo源代码}

\begin{lstlisting}[language=c]
kk=2;
[mdd,ndd]=size(dd);
while ~isempty(V)
    [tmpd,j]=min(W(i,V));tmpj=V(j);
for k=2:ndd
    [tmp1,jj]=min(dd(1,k)+W(dd(2,k),V));
    tmp2=V(jj);tt(k-1,:)=[tmp1,tmp2,jj];
end
    tmp=[tmpd,tmpj,j;tt];[tmp3,tmp4]=min(tmp(:,1));
if tmp3==tmpd, ss(1:2,kk)=[i;tmp(tmp4,2)];
else,tmp5=find(ss(:,tmp4)~=0);tmp6=length(tmp5);
if dd(2,tmp4)==ss(tmp6,tmp4)
    ss(1:tmp6+1,kk)=[ss(tmp5,tmp4);tmp(tmp4,2)];
else, ss(1:3,kk)=[i;dd(2,tmp4);tmp(tmp4,2)];
end;
end
    dd=[dd,[tmp3;tmp(tmp4,2)]];V(tmp(tmp4,3))=[];
    [mdd,ndd]=size(dd);
    kk=kk+1;
end;
S=ss;
D=dd(1,:);
 \end{lstlisting}
\end{appendices}

\end{document} 